\documentclass{leveldoc}
\DocUses{foundations,slingshot,two_stages,injection,stray_field,rf_kick,rf_separator,streak_camera}
\DocFingerprint{2145ee355d7eb0c8}

\begin{document}

\section{What the level asks}\label{sec:asks}

The game's text:
\leveldescription

Two identical ions ($q = \val{q}$, $m = \val{m}$, $T_0 = \val{Tzero}$) leave
$(\val{xZero}, \val{yZero})$ along $+x$ in a world with $c = \val{c}$, the second
$\Delta t = \val{dt}$ after the first: half the period $T = 2\pi/\omega = \val{period}$ of the level's
antenna frequency $\omega = \val{omega}$. The game's ``bunches'' are single ions. Each ion is flown
under both of the level's disturbances (\physref{2.3}), ``\val{distOff}'', no outside field, and
``\val{distOn}'', a uniform field $\vect{E} = (0, \val{Ey})$; the level is solved only if all four
flights succeed. Each flight must first pass the gate $\val{xGLo} \le x \le \val{xGHi}$,
$\val{yGLo} \le y \le \val{yGHi}$, entering it with a direction within $\delta = \val{gHalf}^\circ$
of the $+x$ axis (by the rules \xeqref{two_stages}{eq:gate} and \xeqref{injection}{eq:cone}; an
entry outside the cone does not count, and a later entry is tested again, \physref{6.2}), and then
enter its own
detector: the first ion $\val{xLoOne} \le x \le \val{xHiOne}$, $\val{yLoOne} \le y \le \val{yHiOne}$
above the beam line, the second $\val{xLoTwo} \le x \le \val{xHiTwo}$,
$\val{yLoTwo} \le y \le \val{yHiTwo}$ below it. A detector entered before the gate is a skipped gate.
The player places up to $\val{maxQ}$ charges, of magnitudes $\{\val{magnitudes}\}$ and either sign,
and up to $\val{maxA}$ antennas, as in \levelref{rf_kick} (at the frequency $\omega$, amplitudes
$\{\val{amplitudes}\}$ of either sign, oriented at $0^\circ$, $45^\circ$, $90^\circ$ or $135^\circ$),
at nodes with $\val{xRegLo} \le x \le \val{xRegHi}$, $\val{yRegLo} \le y \le \val{yRegHi}$. A flight
ends when the ion enters its detector, touches a placed element's body (radius $\val{Rq}$ for a
charge, $\val{Ra}$ for an antenna), leaves the world, or reaches $t = \val{tMax}$
(\xsecref{foundations}{sec:board}).

The reference places two charges, $Q_1 = \val{Q1}$ at $(\val{xQ1}, \val{yQ1})$ and
$Q_2 = \val{Q2}$ at $(\val{xQ2}, \val{yQ2})$, and two antennas oriented at $\val{angA1}^\circ$,
$p_1 = \val{pA1}$ at $(\val{xA1}, \val{yA1})$ and $p_2 = \val{pA2}$ at $(\val{xA2}, \val{yA2})$. The
document sizes the stray field against the gate, which shows what the steering must do (centre the
two cases on the gate in height and in angle); proves that no static field can separate the ions
and that the late ion is the early one with the antennas reversed; solves the empty board in closed
form; and marches the four flights.

\section{Setup}\label{sec:setup}

The level file, field by field (the appendix reproduces it):
\setuptable

By \xeqref{foundations}{eq:launch} each ion starts with $\gamma_0 = \val{gammaZero}$, momentum
$p_\mathrm{i} = \val{pZero}$ and speed $v_0 = \val{vZero}$.

\begin{figure}[htbp]
\centering
\begin{tikzpicture}
\begin{axis}[axis equal image, width=0.95\linewidth, xmin=-1.6, xmax=41.6, ymin=-1.6, ymax=21.6,
  xtick={0,5,...,40}, ytick={0,5,...,20}, tick label style={font=\small},
  xlabel={$x$}, ylabel={$y$}, label style={font=\small}]
  \draw[black!8] (0,0) grid (40,20);
  \draw[black!45, dashed] (-1,-1) rectangle (\rawval{xWorld},\rawval{yWorld});
  \draw[black!30, dotted] (\rawval{xRegLo},\rawval{yRegLo}) rectangle (\rawval{xRegHi},\rawval{yRegHi});
  \draw[black!60, thick] (\rawval{xGLo},\rawval{yGLo}) rectangle (\rawval{xGHi},\rawval{yGHi});
  \fill[blue!60!black, fill opacity=0.18] (\rawval{xLoOne},\rawval{yLoOne}) rectangle (\rawval{xHiOne},\rawval{yHiOne});
  \fill[red!60!black, fill opacity=0.18] (\rawval{xLoTwo},\rawval{yLoTwo}) rectangle (\rawval{xHiTwo},\rawval{yHiTwo});
  \addplot[thick, blue!70!black] table {pathoneoff.dat};
  \addplot[thick, blue!70!black, dashed] table {pathoneon.dat};
  \addplot[thick, red!70!black] table {pathtwooff.dat};
  \addplot[thick, red!70!black, dashed] table {pathtwoon.dat};
  \fill[red!80!black] (\rawval{xQ1},\rawval{yQ1}) circle[radius=\rawval{Rq}];
  \fill[blue!80!black] (\rawval{xQ2},\rawval{yQ2}) circle[radius=\rawval{Rq}];
  \fill[orange!80!black] (\rawval{xA1},\rawval{yA1}) circle[radius=\rawval{Ra}];
  \fill[orange!80!black] (\rawval{xA2},\rawval{yA2}) circle[radius=\rawval{Ra}];
  \fill (\rawval{xZero},\rawval{yZero}) circle[radius=0.15];
\end{axis}
\end{tikzpicture}
\caption{The reference: the positive charge $Q_1$ (red disc), the negative $Q_2$ (blue), the two
antennas (orange), the placement region (dotted), the gate (outlined) and the detectors. The first
ion (blue) and the second, $\Delta t$ later (red), with the stray field off (solid) and on
(dashed).}
\label{fig:board}
\end{figure}

\section{Principles and model}\label{sec:model}

\begin{principles}
  \item The charges' fields \xeqref{foundations}{eq:coulomb}, the uniform outside field of each
    disturbance (\physref{2.3}), as in \levelref{stray_field}; the antennas' exact retarded fields,
    \xeqref{rf_kick}{eq:fields}; a flight launched at $t_0$ in the fields of the lab time $t_0 + t$,
    \xeqref{rf_separator}{eq:offset}.
  \item The relativistic equations of motion \xeqref{foundations}{eq:eom} with the launch
    \xeqref{foundations}{eq:launch}; Liénard's power \xeqref{foundations}{eq:larmor}.
  \item The gate \xeqref{two_stages}{eq:gate} and its direction cone \xeqref{injection}{eq:cone},
    tested at every entry (\physref{6.2}).
  \item The board's rules; the classical Runge--Kutta method, cubic Hermite interpolation and the
    error estimate from the step $2h$ (\xsecref{foundations}{sec:board},
    \xsecref{foundations}{sec:numerics}), as in \xsecref{slingshot}{sec:method}.
\end{principles}

\begin{assumption}{The ions' radiation, and fixed sources}\label{as:rf}
As in \levelref{rf_kick}, whose assumptions \xref{rf_kick}{as:radiation} and
\xref{rf_kick}{as:driven} these are: the ions' radiation is neglected, at most $\val{radOneOn}$ here,
and the antennas oscillate at fixed amplitudes. The charges are held in place, the outside field is
uniform and static under each disturbance (an exact solution of Maxwell's equations in vacuum), and
the ions fly separately. These are the game's model.
\end{assumption}

\section{Derivation}\label{sec:derivation}

\subsection{The stray field against the gate}\label{sec:stray}

The supply's field pushes with $F = qE_y = \val{force}$, a transverse acceleration of
$F/(\gamma_0 m) = \val{accel}$ at the start. Along the undeflected line the ion reaches the gate's
face at $t = \val{tGate}$: by then the push alone has lowered it by about
$\tfrac12 (F/\gamma_0 m)t^2 = \val{dropGate}$ and turned it by $\arctan(\abs{F}t/p_\mathrm{i}) = \val{angleGate}^\circ$, against a gate
$\val{gateHeight}$ high and a cone $2\delta = \val{coneWidth}^\circ$ wide. The charges and antennas
act the same under both disturbances, so whatever they do, the stray field changes the entry by
nearly the gate's whole height and the cone's whole width: a solution must enter the gate near its
top, climbing at nearly $+\delta$, with the field off, and near its bottom, descending at nearly
$-\delta$, with it on. The reference's four entries (the march, \S\ref{sec:march}; $\theta$ is the
direction of motion above the axis):
\begin{center}\small\val{gates}\end{center}
Along the reference's paths the field changes the entry by $\val{dYOne}$ in height and
$\val{dThetaOne}^\circ$ in angle for the first ion ($\val{dYTwo}$ and $\val{dThetaTwo}^\circ$ for the
second), a little less than along the straight line; the field-off entries lie within
$\val{gateTop}$ of the gate's top, the field-on ones $\val{gateBottom}$ above its bottom, and the
smallest angular margin is $\val{minGateMargin}^\circ$.

\subsection{Two ions in one static field}\label{sec:late}

The charges' fields and the stray field do not depend on time. With them alone the two ions, which
differ only in their launch time, obey the same equations \xeqref{rf_separator}{eq:offset} with the
same initial state: they fly the same path and cannot reach two different detectors, whatever the
charges. The antennas' moments satisfy $\vect{p}(t + T/2) = -\vect{p}(t)$, and their fields
\xeqref{rf_kick}{eq:fields} are linear in $\vect{p}$ and its derivatives at the retarded time, so the
second ion, launched at $t_0 = T/2$, moves in
\begin{equation}\label{eq:late}
  \vect{E}_Q + \vect{E}_\mathrm{out} + \vect{E}_\mathrm{ant}(\vect{x}, T/2 + t;\, \vect{p}_0)
  = \vect{E}_Q + \vect{E}_\mathrm{out} + \vect{E}_\mathrm{ant}(\vect{x}, t;\, -\vect{p}_0),
\end{equation}
and the same for $\vect{B}$: it flies as an ion launched with the first one would, with every
antenna's amplitude reversed. Unlike the second ion of \levelref{rf_separator}, which \xeqref{rf_separator}{eq:conjugate}
makes the opposite charge, it is not the opposite charge, for the static fields keep their sign. Marching both, with the second ion's detector, gives the same end state to
$\val{lateGap}$. Only the antennas can send the two ions apart, and they must do it while the
static fields keep both on a path through the gate.

\subsection{The empty board in closed form}\label{sec:hyperbolic}

With no element and the field off the ions fly straight along $y = \val{yZero}$, through the gate
on its axis, to the world's right edge at $t = \val{bareTOff}$. With the field on, the force
$(0, F)$ is constant: $p_x = p_\mathrm{i}$, $p_y = Ft$, and the energy
$\varepsilon = \sqrt{\varepsilon_0^2 + c^2F^2t^2}$ with $\varepsilon_0 = \gamma_0 mc^2$. Integrating
$\dot{\vect{x}} = c^2\vect{p}/\varepsilon$,
\begin{equation}\label{eq:hyperbolic}
  x - x_0 = \frac{p_\mathrm{i}c}{\abs{F}}\operatorname{arsinh}\frac{c\abs{F}t}{\varepsilon_0},
  \qquad
  y - y_0 = \frac{\varepsilon - \varepsilon_0}{F},
\end{equation}
the second also from the energy balance $\dd\varepsilon = F\,\dd y$. The ion passes below the gate
(\S\ref{sec:stray}) and reaches the bottom edge $y = \val{bareYOn}$ when
$\varepsilon = \varepsilon_0 + F(\val{bareYOn} - y_0)$: at $t = \val{bareTOn}$, at
$x = \val{bareXOn}$, just before the right edge.

\subsection{The march}\label{sec:march}

Each flight is marched by \xeqref{foundations}{eq:rk4} at $h = \val{H}$ with Liénard's energy as one
more state. Every entry into the gate and the detector is located on the cubic Hermite interpolant by
bisection; at a gate entry the momentum, from a Runge--Kutta step of the located length, is tested
against the cone, and the flight goes on. The error is estimated against the step $2h$, as in
\xsecref{rf_kick}{sec:flight}.

\section{The reference solution by hand}\label{sec:hand}

\begin{procedure}
  \item\label{st:stray} \emph{The stray field's size.} \S\ref{sec:stray}: the field moves the entry
    by nearly the gate's height and the cone's width.
  \item\label{st:gate} \emph{The gate.} The march: the four entries of the table in
    \S\ref{sec:stray}, each within the cone at its first entry.
  \item\label{st:oneoff} \emph{The first ion, field off.} It enters through the \val{faceOneOff}
    face of its detector at $(\val{xOneOff}, \val{yOneOff})$ at $t = \val{tOneOff}$, with
    $\vect{p} = (\val{pxOneOff}, \val{pyOneOff})$ (step estimate $\val{errOneOff}$).
  \item\label{st:oneon} \emph{The first ion, field on.} Through the \val{faceOneOn} face at
    $(\val{xOneOn}, \val{yOneOn})$ at $t = \val{tOneOn}$, with $\vect{p} = (\val{pxOneOn}, \val{pyOneOn})$
    (step estimate $\val{errOneOn}$).
  \item\label{st:twooff} \emph{The second ion, field off.} Launched $\Delta t$ later: through the
    \val{faceTwoOff} face at $(\val{xTwoOff}, \val{yTwoOff})$ at $t = \val{tTwoOff}$ after its
    launch, with $\vect{p} = (\val{pxTwoOff}, \val{pyTwoOff})$ (step estimate $\val{errTwoOff}$).
  \item\label{st:twoon} \emph{The second ion, field on.} Through the \val{faceTwoOn} face at
    $(\val{xTwoOn}, \val{yTwoOn})$ at $t = \val{tTwoOn}$ after its launch, with
    $\vect{p} = (\val{pxTwoOn}, \val{pyTwoOn})$ (step estimate $\val{errTwoOn}$).
\end{procedure}

\begin{result}[The reference flights]
Under both disturbances each ion passes the gate within the cone and enters its own detector. The
level is solved.
\end{result}

\section{Comparison with the game}\label{sec:comparison}

\comparisontable

All trajectory differences on the reference are at most $\val{largest}$, in the first ion's flight
with the field on, which passes about $\val{closeQ}$ from the negative charge, where the force is
strongest and both integrations are least accurate: there our own step estimate is
$\val{errOneOn}$, and the difference is of its size. The tolerances are ten times the larger of the
two computations' estimates, or the game's own preview--verification gap where that is larger; the
radiated energies agree to about $10^{-9}$ relative. The empty board agrees with
\eqref{eq:hyperbolic} to the game's accuracy. With the reference's charges alone both ions pass the
gate under both disturbances and fly the same path, as \S\ref{sec:late} says, leaving the world at
height $\val{chargesYOff}$ with the field off and $\val{chargesYOn}$ with it on: the game's two shots
give the same numbers.

\section{Other solutions}\label{sec:others}

The level separates into its two stages only roughly, for the antennas also act before the gate and
the charges after it, but \S\ref{sec:stray} and \S\ref{sec:late} hold for every solution: the
static elements must centre the field-off and field-on entries on the gate in height and in angle,
and at least one antenna must kick the two ions apart, the second with the opposite of the first's
antenna fields. The reference does both with two of the five charges and two of the three antennas,
with small margins: $\val{minGateMargin}^\circ$ in angle and $\val{gateTop}$ in height at the gate.
A map of the solutions would take the game's verdicts, which this document does not reproduce.

\section{What the game computes}\label{sec:engine}

The game adds the charges' Coulomb fields, the disturbance's uniform field (\physref{2.3}) and each
antenna's retarded fields at the level's frequency (\physref{2.4}), at the lab time $t_0 + t$, and
integrates the motion with its adaptive eighth-order Runge--Kutta method (\physref{3}, \physref{5},
\physref{6}, \physref{7}); the gate is tested at every entry (\physref{6.2}), and each shot is flown
and verified under each disturbance. The problem is the same as here, and the differences in the
comparison are the two integrations' errors.

\section{Further reading}\label{sec:reading}

One setup under two outside fields: \levelref{stray_field}; gates: \levelref{two_stages}; the
direction cone: \levelref{injection}; the antennas' fields and kicks: \levelref{rf_kick}; two ions
half a period apart: \levelref{rf_separator}; three a third of a period apart:
\levelref{streak_camera}. Relativistic motion in static uniform fields: \Jackson{Ch.~12}.

\levelfileappendix

\end{document}
